\section{Introduction}
Company memory should provide the context a decision needs: the right evidence, at the right time, with its authority and provenance intact. In business-to-business (B2B) marketing, that context includes an approved target market, buyer roles, product claims, proof and observations about prospective accounts. A domain-specialized memory may organize these relationships explicitly; a general-purpose company memory may serve marketing alongside other functions. Whether specialization improves context provision is a different question from whether a model recommends a useful action from that context.

This preliminary component study uses a small shared catalog of research actions as an auditable test instrument. Each action has explicit prerequisites and a resource allowance. A reader proposes one supported action or requests missing evidence; a validator checks its sources and can dry-run its steps against fictional providers. The catalog formalizes what context is needed. It does not define the commercially best next action. Historical prompts and machine-readable traces call these workflows \emph{Plays}; that identifier is retained in frozen artifacts for exact reproduction, not required by the research question.

The experiment is a standalone prototype with fictional records, motivated by marketing decision support. It neither evaluates a production product nor establishes that generic company memory lacks domain semantics. The mechanism tested here selects complementary evidence for declared prerequisites; its implementation does not require a graph database.

The motivating question is whether task-specific memory structure improves this interface beyond capable generic retrieval and a domain prompt. The question is easy to confound. A specialized system may receive cleaner facts, a larger context, an extra model call, or a better planner. A generic system may be handicapped by missing domain vocabulary. High retrieval recall can also coexist with incorrect recommendations. We therefore hold the verified facts, metadata boundary, catalog, reader and output interface constant, and compare retrieval selection under an explicit total-input allowance. We include a prompt-only specialization control, a domain-neutral diversity control, and separately priced full context.

The contribution is a controlled empirical study and an executable fixture suite. Marketing memory and planning, domain-aware memory, belief-based recommendation, complementary evidence retrieval and derived-state maintenance all have close precedents. Our narrower question concerns evidence coverage under changing policy, observations and permissions. The selection mechanism can operate on any catalog of typed prerequisites, so a benefit would not establish a uniquely marketing algorithm or an ontology-based commercial advantage.

We distinguish a \emph{system of record}, which preserves evidence and approved policy; \emph{context provision}, which selects sufficient current support; and \emph{recommendation}, which uses that support to choose an action. A downstream executable check verifies a formal specification. The existing frozen study records context availability and a shared reader's decisions; its original primary outcome remains decision correctness. The revised presentation does not retroactively turn context recall into a preregistered primary endpoint. Follow-up memory research should score context directly with the reader fixed, while recommendation research should vary the decision method with context fixed. Neither the present component study nor that separation establishes strategic usefulness or campaign outcomes.
